\chapter{能量、曲率与辐射}

\lectureoverview{一旦去掉两个限制，牛顿模型就会变得更有启发性。第一，共动星系的总能量不必恰好为零；让其能量取任意值会引入 \(a^{-2}\) 项，广义相对论将把它重新解释为空间曲率。第二，宇宙源不必是无压物质。光子波长随尺度因子拉伸，因此辐射按 \(a^{-4}\) 稀释，并给出规律 \(a(t)\propto t^{1/2}\)。因此，同时包含物质与辐射的宇宙在小 \(a\) 时由辐射主导，在较大的 \(a\) 时由物质主导。在此过程中，我们将把局部牛顿运动与整体相对论几何区分开来，并找出那些没有广义相对论就无法解决的问题。}

\section{牛顿能看到的}

在扩展模型之前，我们应当问：为什么先前的牛顿计算竟然有效？物理宇宙由广义相对论描述，其空间几何不必全局为欧几里得几何。然而，每一个光滑弯曲空间在足够小的邻域内都近似平坦。地球表面的一小块看起来是平面的；同样，远小于其曲率半径的宇宙区域也可以当作平坦来处理。

``小''在这里是相对于曲率尺度而言的。这一小块区域仍可能横跨数十亿光年。在其中，当相邻星系的相对本动速度为非相对论性的、且引力场很弱时，就可以用牛顿力学处理它们。因此，牛顿计算并不是爱因斯坦理论的对手。它是更完整描述的局部、慢速极限。

这里立刻有一个警告。辐射并不缓慢运动。光子，以及仍处于相对论状态的中微子，不能被当作普通牛顿粒子的气体来处理。牛顿式推理仍将继续组织膨胀方程，但其右侧的源必须从静质量密度扩大到能量密度。正是在这里，近似开始暴露出自身的边界。

\section{回顾尺度因子}

取一个均匀分布的星系系统，并把坐标附着在一个随宇宙平均流动而运动的网格上。位于固定共动坐标 \(x\) 处的星系，到所选原点的物理距离为
\begin{equation}
D(t)=a(t)x
\label{eq:c02-distance}
\end{equation}
尺度因子 \(a(t)\) 把坐标间隔转换为物理间隔。

\(a\) 的绝对值是约定性的。如果所有共动坐标都按一个常数重新标度，那么同样的物理距离可以由 \(a\) 的逆标度来表示。因此，\(a\) 和赋予一个坐标格胞的质量都不是各自不变的。尺度因子的比值以及分数膨胀率是不变的：
\begin{equation}
\frac{a(t_2)}{a(t_1)},
\qquad
H(t)\equiv\frac{\dot a}{a}.
\label{eq:c02-hubble}
\end{equation}
确实，如果 \(a'=\alpha a\)，其中 \(\alpha\) 为固定常数，那么 \(\dot a'/a'=\dot a/a\)。

设 \(\nu_m\) 为单位共动坐标体积内包含的质量。如果重新标记网格，它的数值会改变；但当普通物质守恒时，它不随时间改变。一个共动单位立方体的物理体积为 \(a^3\)，因此物质密度为
\begin{equation}
\boxed{\rho_m(a)=\frac{\nu_m}{a^3}.}
\label{eq:c02-matter-density}
\end{equation}
这个 \(a^{-3}\) 稀释不过是体积增长：共动盒中的静能量保持固定，而它的体积在增长。

真实星系并不是被刚性钉在网格上的。它们在群中轨道运动，落入星系团，并相对于平均流动带有小速度。尺度因子描述的是大尺度上的平滑运动；这种局部偏离称为 \emph{本动速度}。

\begin{classroomqa}
\audiencequestion{每个星系都必须恰好沿径向远离其他每个星系吗？}
\lecturerresponse{只有平均哈勃流才具有这种形式。一个星系的实际运动是共同膨胀与由附近结构引起的本动速度之和。这种区别就像河流水流上叠加了涡旋：涡旋在局部很重要，但不会抹去大尺度流动。\lecturetimestamp{00:24:47--00:26:33}}
\end{classroomqa}

对于牛顿总能量为零的无压物质，上一讲得到
\begin{equation}
H^2=\frac{8\pi G}{3}\rho_m
=\frac{8\pi G\nu_m}{3a^3}.
\label{eq:c02-zero-energy-friedmann}
\end{equation}
由于整体网格归一化是任意的，我们可以暂时选择使 \(8\pi G\nu_m/3=1\) 的单位。在膨胀分支上，
\begin{equation}
\frac{\dot a}{a}=\frac{1}{a^{3/2}},
\qquad
\dot a=\frac{1}{\sqrt a}.
\end{equation}
将导数反演并积分，得到
\begin{equation}
\frac{dt}{da}=\sqrt a,
\qquad
t-t_B=\frac{2}{3}a^{3/2},
\end{equation}
其中 \(t_B\) 是这个理想化解达到 \(a=0\) 的时间。因此
\begin{equation}
\boxed{a(t)\propto (t-t_B)^{2/3}.}
\label{eq:c02-matter-law}
\end{equation}
在这个仅含物质的临界模型中，宇宙永远膨胀。它的膨胀会减慢，因为 \(\dot a\propto a^{-1/2}\)，但 \(\dot a\) 不会在任何有限的 \(a\) 值处变为零。

\begin{figure}[ht]
\centering
\includegraphics[width=0.68\textwidth]{lecture_02_figure_02.png}
\caption{复习板上记录了\(\dot a=1/\sqrt a\)以及框出的物质主导定律\(a\propto t^{2/3}\)。中间的代数推导部分被遮挡，由方程~\eqref{eq:c02-zero-energy-friedmann}--\eqref{eq:c02-matter-law}补充。}
\label{fig:c02-matter-recap}
\lectureframe{00:29:20.838}
\end{figure}

\begin{classroomqa}
\audiencequestion{为什么要舍弃负的平方根？}
\lecturerresponse{方程并没有舍弃它。方程~\eqref{eq:c02-zero-energy-friedmann}确定的是\(H^2\)，而不是\(H\)的符号。正根描述膨胀，负根描述收缩。行星引力场中的一块岩石可以向外运动，也可以向内运动，同时满足同一个能量方程；它的初始条件选择了分支。我们选择了膨胀分支，因为那正是所讨论的宇宙。\lecturetimestamp{00:23:09--00:24:46}}
\end{classroomqa}

下一个模型仍将忽略暗能量。这种忽略是有意为之：我们先理解较早的物质与辐射框架，然后再加入控制当前加速时期的成分。相比之下，暗物质已经包含在物质项中，因为非相对论性暗物质也按\(a^{-3}\)稀释。

\begin{classroomqa}
\audiencequestion{暗能量难道不是今天的主导成分吗？}
\lecturerresponse{是的。这里展开的双成分叙述是发现当前加速之前的标准教学宇宙学。暗物质属于物质；暗能量是第三个成分，将单独引入。从较旧的模型开始，可以让我们看到每增加一项会改变什么。\lecturetimestamp{00:28:02--00:28:45}}
\end{classroomqa}

\section{恢复能量常数}

零能量选择是一种特殊的初始条件。为了去掉它，回到均匀球体边界上一个质量为\(m\)的星系。设它的共动半径为\(x\)，于是
\begin{equation}
D=ax,
\qquad
\dot D=\dot a x.
\label{eq:c02-comoving-particle}
\end{equation}
牛顿壳层定理表明，只有球体内部的质量\(M(D)\)影响径向加速度；外部球壳不施加净力。该星系守恒的牛顿能量为
\begin{equation}
\frac{1}{2}m\dot D^{\,2}-\frac{GM(D)m}{D}=E.
\label{eq:c02-newton-energy}
\end{equation}
因子\(m\)在两个项中都应出现。在板上快速书写时很容易漏掉它，但保留它就能清楚地表明，最终的常数是单位检验质量的能量。

均匀性允许我们把代表性星系选在\(x=1\)处。于是\(D=a\)，\(\dot D=\dot a\)，并且该球体有
\begin{equation}
V_{\rm sph}=\frac{4\pi}{3}a^3,
\qquad
M(a)=\rho_m V_{\rm sph}=\frac{4\pi}{3}\rho_m a^3.
\label{eq:c02-enclosed-mass}
\end{equation}
板书把\(V\)既用于体积，又在一个旁注中用于径向速度。我们将保留无歧义的符号\(V_{\rm sph}\)和\(\dot D\)。

\begin{figure}[ht]
\centering
\includegraphics[width=0.86\textwidth]{lecture_02_figure_03.png}
\caption{包含质量的球体与非零能量的弗里德曼方程。可见的标记包括\(V_{\rm sph}=4\pi a^3/3\)、中心质量\(M\)、\(x=1\)、\(D=a\)和\(\dot D=\dot a\)。方程的一部分被讲师遮住，下面用代数方法重建。}
\label{fig:c02-energy-board}
\lectureframe{00:37:59.876}
\end{figure}

\begin{figure}[ht]
\centering
\begin{tikzpicture}[scale=1.05,>=Latex]
  \draw[thick] (0,0) circle (2.6);
  \fill (0,0) circle (2.2pt) node[below left=2pt] {\(M(a)\)};
  \draw[thick,->] (0,0) -- (2.58,0);
  \fill (2.6,0) circle (2.2pt) node[right=4pt] {galaxy \(m\)};
  \node[above] at (1.30,0.04) {\(D=ax\)};
  \node[below] at (1.30,-0.04) {\(x=1\Rightarrow D=a\)};
  \draw[->,blue!60!black,thick] (2.63,0.18) -- (3.55,0.72)
    node[right] {\(\dot D=\dot a\)};
  \node at (0,-3.05) {\(M(a)=\frac{4\pi}{3}\rho_m a^3\)};
\end{tikzpicture}
\caption{在\(x=1\)处对牛顿球体的清晰重建。外部球壳相互抵消，因此边界星系只响应内部包含的质量。}
\label{fig:c02-sphere-redraw}
\end{figure}

将\(D=a\)代入方程~\eqref{eq:c02-newton-energy}，除以\(m/2\)，并定义常数
\begin{equation}
C\equiv\frac{2E}{m}.
\end{equation}
我们得到
\begin{equation}
\dot a^{\,2}-\frac{2GM(a)}{a}=C.
\label{eq:c02-energy-a}
\end{equation}
除以\(a^2\)并利用方程~\eqref{eq:c02-enclosed-mass}，得到
\begin{equation}
\left(\frac{\dot a}{a}\right)^2-
\frac{8\pi G}{3}\rho_m=\frac{C}{a^2},
\label{eq:c02-energy-form}
\end{equation}
或者等价地，
\begin{equation}
\boxed{H^2=\frac{8\pi G}{3}\rho_m+\frac{C}{a^2}.}
\label{eq:c02-friedmann-C}
\end{equation}
这就是无压物质的第一个弗里德曼方程，以牛顿能量约定写出。新的\(a^{-2}\)项记录了代表性星系是以高于、等于还是低于逃逸能量的方式发射的。

代入\(\rho_m=\nu_m/a^3\)后，方便地定义
\begin{equation}
A_m\equiv\frac{8\pi G\nu_m}{3}>0,
\end{equation}
于是
\begin{equation}
H^2=\frac{A_m}{a^3}+\frac{C}{a^2},
\qquad
\dot a^{\,2}=\frac{A_m}{a}+C.
\label{eq:c02-energy-branches}
\end{equation}
与其立即积分，不如按极限来读这个方程，它更有启发性。

\subsection{正能量}

如果\(C>0\)，则在\(a\)很小时，物质项\(A_m/a^3\)占主导。因此早期膨胀保持物质定律\(a\propto t^{2/3}\)。然而在\(a\)很大时，\(C/a^2\)项在\(H^2\)中占主导，或者等价地
\begin{equation}
\dot a^{\,2}\longrightarrow C.
\end{equation}
膨胀趋近于恒定速率，
\begin{equation}
a(t)\sim\sqrt C\,t.
\label{eq:c02-positive-late}
\end{equation}
这是以超过逃逸速度抛出一个苹果的宇宙学对应。引力不断减小它的速率，但远处苹果仍保留非零的渐近速度。

\subsection{零能量}

对于\(C=0\)，方程~\eqref{eq:c02-energy-branches}回到临界解\(a\propto t^{2/3}\)。它无界膨胀，而\(\dot a\)趋于零。这正是逃逸能量的情形。

\subsection{负能量}

如果\(C=-|C|<0\)，那么
\begin{equation}
H^2=\frac{A_m}{a^3}-\frac{|C|}{a^2}
=\frac{A_m-|C|a}{a^3}.
\label{eq:c02-negative-energy}
\end{equation}
标度因子不能超过
\begin{equation}
a_{\max}=\frac{A_m}{|C|},
\label{eq:c02-amax}
\end{equation}
因为右端在那里变为零。膨胀分支到达\(a_{\max}\)时\(\dot a=0\)；其延续是收缩分支。在牛顿图像中，这是一个以低于逃逸速度抛出的苹果：它上升、转向并落下。

\begin{classroomqa}
\audiencequestion{当右侧等于\(H^2\)、因而不能为负时，模型如何能够重新坍缩？}
\lecturerresponse{它永远不会变为负值。在膨胀过程中，右侧在\(a_{\max}\)处减小到零。那里\(H=0\)。于是解从正平方根变为负平方根，并在\(\dot a<0\)的情况下沿允许区域回溯。一个诱人但不正确的延续会让该表达式穿过零。它不能：零是边界，而不是一个实的\(H\)能够继续穿过的值。\lecturetimestamp{00:48:43--00:53:03}}
\end{classroomqa}

因此，三种牛顿能量给出三种定性历史：\(C<0\)对应重新坍缩，\(C=0\)对应临界逃逸，\(C>0\)对应\(\dot a\)渐近为常数的永恒膨胀。令人惊讶的下一步是，广义相对论为这些相同情形赋予了几何意义。

\section{从能量到曲率}

把方程~\eqref{eq:c02-friedmann-C}中的能量项移到几何一侧，并定义
\begin{equation}
k\equiv-C.
\label{eq:c02-k-definition}
\end{equation}
于是
\begin{equation}
\boxed{H^2+\frac{k}{a^2}=\frac{8\pi G}{3}\rho.}
\label{eq:c02-friedmann-k}
\end{equation}
黑板上在\(C\)、\(K\)和\(\kappa\)之间交替，并且在把该项移过等号时吸收了一个符号。方程~\eqref{eq:c02-k-definition}确定了我们的约定：正的牛顿能量对应\(k<0\)，零能量对应\(k=0\)，负的牛顿能量对应\(k>0\)。

\begin{figure}[ht]
\centering
\includegraphics[width=0.76\textwidth]{lecture_02_figure_04.png}
\caption{将弗里德曼方程整理为左侧几何、右侧源。该画面支持这种概念划分；它并未展示完整的张量爱因斯坦方程。}
\label{fig:c02-geometry-source}
\lectureframe{00:55:22.810}
\end{figure}

这种安排是爱因斯坦方程的一个缩微版本。左侧描述几何：通过\(H\)描述膨胀，通过\(k/a^2\)描述空间曲率。右侧描述物质源。在牛顿力学中，该源是质量密度。在相对论中，它是能量和动量，由应力-能量张量完整表示。对于均匀且各向同性的宇宙，第一弗里德曼方程依赖于总能量密度\(\rho\)。

如果保留普通单位，质量密度通过其静止能量\(\rho_{\rm mass}c^2\)作出贡献。从这里开始，我们设\(c=1\)，因此质量和能量使用相同单位。符号\(\rho\)于是可以包括静止质量、动能、辐射以及任何其他均匀能量组分。

\begin{figure}[ht]
\centering
\includegraphics[width=0.88\textwidth]{lecture_02_figure_05.png}
\caption{弗里德曼方程的并排黑板形式：使用\(\rho\)的曲率形式，以及经过\(\rho_m=\nu_m/a^3\)后得到的物质形式。黑板上的记号发生变化；方程~\eqref{eq:c02-k-definition}在此保持符号约定固定。}
\label{fig:c02-curvature-comparison}
\lectureframe{00:57:52.659}
\end{figure}

几何解释是全局性的。在相对论模型中，\(k=0\)描述平坦空间切片，\(k>0\)描述正曲率球面切片，\(k<0\)描述负曲率双曲切片。牛顿推导再现了相同的演化方程，却不具有这些全局几何。因此，它的非零能量常数极其精确地模仿了曲率，但并不是对弯曲空间本身的推导。

\begin{classroomqa}
\audiencequestion{如果牛顿力学和广义相对论的几何不同，它们如何能给出相同的膨胀方程？}
\lecturerresponse{广义相对论包含平坦情形，而且局部地看，每一种光滑几何都近似平坦。更令人惊讶的是，牛顿能量常数进入演化方程的位置，恰好就是相对论空间曲率所占据的位置。动力学公式一致，而全局解释并不一致：欧几里得牛顿空间并没有暗中变成球面或双曲空间。这三种几何将不得不在广义相对论内部构造出来。\lecturetimestamp{01:20:59--01:24:20}}
\end{classroomqa}

\section{膨胀盒子中的辐射}

现在我们改变宇宙的内容物。考虑一个共动单位立方体。它的物理体积是\(a^3\)。对于非相对论物质，该盒子中的粒子数固定，且每个粒子的静止能量近似固定，从而给出\(\rho_m\propto a^{-3}\)。对于辐射，还有另一种效应：每个光子随着宇宙膨胀而损失能量。

波长为\(\lambda\)的光子具有能量
\begin{equation}
E_\gamma=\frac{hc}{\lambda},
\qquad
E_\gamma\propto\lambda^{-1}.
\label{eq:c02-photon-energy}
\end{equation}
在均匀膨胀下，自由传播的波长随尺度因子拉伸：
\begin{equation}
\frac{\lambda(t_2)}{\lambda(t_1)}
=\frac{a(t_2)}{a(t_1)},
\qquad
\lambda\propto a.
\label{eq:c02-redshift}
\end{equation}
这就是宇宙学红移。它意味着
\begin{equation}
E_\gamma\propto a^{-1}.
\label{eq:c02-photon-redshift}
\end{equation}
如果共动盒中的光子数守恒，则其总辐射能按\(a^{-1}\)下降。除以其物理体积\(a^3\)得到
\begin{equation}
\boxed{\rho_r(a)=\frac{\nu_r}{a^4}.}
\label{eq:c02-radiation-density}
\end{equation}
这里\(\nu_r\)是自由流动辐射密度的恒定振幅。辐射比物质多一个\(a\)的负幂次：三个幂次来自空间稀释，一个来自每个光子的红移。

吉他弦的类比给出了有用的直觉。拟合到弦上的模式具有整数个半波长。如果弦被缓慢拉长，该整数不能连续变化；相反，波长会调整以适应新的长度。同样，对于缓慢膨胀的环上的波，节点数是绝热不变量，因此其波长跟随周长。一般的波可以分解为这样的模式。

这并不能替代在膨胀时空中求解麦克斯韦方程组。它是方程~\eqref{eq:c02-redshift}背后受控的直觉：缓慢的几何膨胀改变物理波长，同时保持模式标签不变。

\begin{classroomqa}
\audiencequestion{这个论证需要一个导电或反射的盒子吗？在没有物理壁的空膨胀空间中会发生什么？}
\lecturerresponse{壁是一种记账工具，而不是声称宇宙被导体包围。完整的证明来自膨胀度规中的波传播。对于目前的论证，我们使用观测到的以及相对论推导出的事实：自由传播的波长随\(a\)缩放。盒子有助于计数光子和物理体积；它不会引起宇宙学红移。\lecturetimestamp{01:02:21--01:04:37}}
\end{classroomqa}

\begin{classroomqa}
\audiencequestion{为什么波长应该跟随网格，而不是保持固定的物理长度？}
\lecturerresponse{考虑一个环上的模式，其周长绝热变化。它的整数节点数不能连续漂移，因此允许的波长必须随环变化。将一般的波展开为傅里叶模式，逐个模式给出相同的结论。这提供了所需的绝热直觉；精确的陈述属于弯曲时空中的场论。\lecturetimestamp{01:34:18--01:38:33}}
\end{classroomqa}

\subsection{辐射主导解}

暂时忽略曲率，将方程~\eqref{eq:c02-radiation-density}代入弗里德曼方程：
\begin{equation}
H^2=\frac{8\pi G\nu_r}{3a^4}
\equiv\frac{A_r}{a^4},
\qquad A_r>0.
\label{eq:c02-radiation-friedmann}
\end{equation}
在膨胀分支上，
\begin{equation}
\frac{\dot a}{a}=\frac{\sqrt{A_r}}{a^2},
\qquad
\dot a=\frac{\sqrt{A_r}}{a}.
\end{equation}
因此
\begin{equation}
a\,da=\sqrt{A_r}\,dt,
\qquad
\frac{a^2}{2}=\sqrt{A_r}(t-t_B),
\end{equation}
并且
\begin{equation}
\boxed{a(t)\propto(t-t_B)^{1/2}.}
\label{eq:c02-radiation-law}
\end{equation}
指数\(1/2\)和\(2/3\)之间的区别在视觉上不大，但在物理上是决定性的。早期宇宙的膨胀率、温度、反应冻结和视界大小都取决于哪个组分占主导。

\section{物质与辐射的宇宙}

对于分别守恒的物质和辐射，它们的密度相加：
\begin{equation}
\rho(a)=\rho_m(a)+\rho_r(a).
\end{equation}
在此比较中忽略曲率，
\begin{equation}
\boxed{H^2=\frac{C_m}{a^3}+\frac{C_r}{a^4},}
\label{eq:c02-mixed-friedmann}
\end{equation}
其中\(C_m,C_r>0\)。板书首先将两个系数标记为\(C_1\)和\(C\)；第二个符号被口头更正为一个不同的常数。方程~\eqref{eq:c02-mixed-friedmann}中的记号使该更正明确。

\begin{figure}[ht]
\centering
\includegraphics[width=0.74\textwidth]{lecture_02_figure_06.png}
\caption{黑板上混合的物质与辐射方程。写下的第二个系数随后被更正为辐射特有的常数，在方程~\eqref{eq:c02-mixed-friedmann}中记为\(C_r\)。}
\label{fig:c02-mixed-board}
\lectureframe{01:11:37.997}
\end{figure}

在小的\(a\)下，辐射项胜出，因为\(a^{-4}\)比\(a^{-3}\)增长得更快：
\begin{equation}
a\ll a_{\rm eq}
\quad\Longrightarrow\quad
H^2\simeq\frac{C_r}{a^4},
\qquad
a(t)\propto t^{1/2}.
\end{equation}
在较大的\(a\)下，物质胜出：
\begin{equation}
a\gg a_{\rm eq}
\quad\Longrightarrow\quad
H^2\simeq\frac{C_m}{a^3},
\qquad
a(t)\propto t^{2/3}.
\end{equation}
相等尺度直接通过令两项相等得到：
\begin{equation}
a_{\rm eq}=\frac{C_r}{C_m}.
\label{eq:c02-equality-scale}
\end{equation}
这种相等与光子退耦不是同一事件。相等问的是两种能量密度何时匹配；退耦问的是相互作用何时停止使物质和辐射保持热接触。

辐射部分包括光子以及任何以超相对论速度运动的物种。非常轻的中微子在足够高的温度下表现如此，而引力子背景会以相同方式缩放。一旦有质量物种变为非相对论性，其主导的静质量贡献就开始像物质一样缩放。冷暗物质属于\(C_m/a^3\)，而不是辐射项。

\begin{classroomqa}
\classquestion{方程~\eqref{eq:c02-mixed-friedmann}中哪一项先占主导，哪一项后占主导？}
\lecturerresponse{对于小的\(a\)，\(a^{-4}\)超过\(a^{-3}\)，因此辐射控制早期膨胀并给出\(a\propto t^{1/2}\)。随着\(a\)增长，辐射与物质之比按\(1/a\)下降；然后物质控制膨胀并给出\(a\propto t^{2/3}\)。\lecturetimestamp{01:11:46--01:13:38}}
\end{classroomqa}

\begin{classroomqa}
\audiencequestion{物质和辐射的相对数量是否被假定为常数？}
\lecturerresponse{它们的密度不是常数，但一旦各组分不再有效地交换能量，且粒子产生或湮灭可以忽略，系数\(C_m\)和\(C_r\)就可以视为常数。此时物质携带守恒的共动静止质量，自由辐射携带守恒的共动光子数，而红移提供了其额外的\(a^{-1}\)因子。在强相互作用或粒子种类湮灭的时期，这种记账方式必须加以修正。\lecturetimestamp{01:13:39--01:14:53}}
\end{classroomqa}

\begin{classroomqa}
\audiencequestion{物质和辐射何时停止强烈相互作用？}
\lecturerresponse{一个粗略的数量级答案是\(10^5\)年，光子与普通物质退耦。精确的复合历史是后续的计算，不应将其与物质--辐射相等混淆。这里所需的要点仅仅是，在退耦之后，这两个组分可以近似地分别追踪。\lecturetimestamp{01:27:29--01:27:46}}
\end{classroomqa}

\begin{classroomqa}
\audiencequestion{我们如何知道在存在星系来示踪膨胀之前，曾发生过一个早期\(a\propto t^{1/2}\)时代？}
\lecturerresponse{我们从它留下的物理记录推断膨胀历史。原初核合成依赖于反应速率与\(H(t)\)之间的早期竞争；微波背景记录了较晚的热状态和扰动。星系观测使物质时代的行为更为直接，而辐射时代则从这些早期宇宙证据中被重构出来。\lecturetimestamp{01:27:46--01:30:10}}
\end{classroomqa}

物质加辐射模型是刻意不完整的。它是宇宙从辐射主导过渡到物质主导的经典框架。当前的加速膨胀需要另一种来源，即暗能量，它将添加一个性质不同的项。

\section{红移能量去向何处}

减小\(E_\gamma\propto a^{-1}\)提出了一个不可避免的问题：光子的能量去了哪里？在一个字面意义上的、带有反射壁的膨胀盒子中，辐射施加压力并对壁做功。热力学陈述是
\begin{equation}
dE=-p\,dV.
\label{eq:c02-pdv}
\end{equation}
随着\(V\)增大，辐射能量减少。损失的能量可能以动能、应变或热的形式出现在任何使壁移动的东西中。

共动盒子没有物质壁。它是一个随平均膨胀运动的假想边界。光子穿过它，但在均匀辐射浴中，进入的统计通量平衡了离开的通量。即使盒子不是容器，它仍然是一个有效的记账体积。压力做功定律仍然描述了共动体积中的能量如何演化。

对于辐射，\(p=\rho/3\)，且\(E=\rho V\)。方程~\eqref{eq:c02-pdv}于是给出
\begin{align}
d(\rho V)&=-\frac{\rho}{3}\,dV,\\
V\,d\rho&=-\frac{4}{3}\rho\,dV,\\
\frac{d\rho}{\rho}&=-\frac{4}{3}\frac{dV}{V}.
\end{align}
由于\(V\propto a^3\)，积分得到\(\rho\propto a^{-4}\)，与波长论证一致。这一热力学推导还清楚地表明，每个光子的能量，而不仅仅是数密度，在减少。

\begin{classroomqa}
\audiencequestion{红移光子的能量去了哪里？宇宙总能量守恒吗？}
\lecturerresponse{有物理壁时，辐射对它们做功，因此这种转移是普通的机械记账。对于整个宇宙，问题更为微妙。全局能量守恒通常来自时间平移对称性，但膨胀的时空一般不具有时间平移不变性，也不必然拥有唯一的守恒总能量。人们可以构建有用的牛顿类比，涉及膨胀的能量，但明确的全局陈述需要广义相对论。两种描述在这里都有用，但从牛顿模型不能得出最终的全局结论。\lecturetimestamp{01:16:51--01:20:58}}
\end{classroomqa}

\begin{classroomqa}
\audiencequestion{如果共动盒子是假想的，且光子穿过其边界，压力做功如何仍然是有用的图像？}
\lecturerresponse{均匀性提供了缺失的记账。平均而言，离开一个共动体积的每个光子，在统计上都被从相邻体积进入的一个光子所替代。局部通量在平均意义上相互抵消，而物理体积和波长发生变化。反射盒子使做功机制具体化；共动盒子在没有字面墙壁的情况下承载相同的平均守恒定律。\lecturetimestamp{01:24:22--01:27:28}}
\end{classroomqa}

\begin{classroomqa}
\audiencequestion{膨胀仅仅是降低了光子密度，还是也降低了每个光子的能量？}
\lecturerresponse{两者都发生。每单位物理体积内的数目按\(a^{-3}\)下降，而宇宙学红移又使每个光子的能量再降低一个因子\(a^{-1}\)。这是绝热冷却的辐射类比。普通气体在缓慢膨胀过程中做功时也会冷却，尽管突然的自由膨胀是另一种热力学过程。\lecturetimestamp{01:42:30--01:45:55}}
\end{classroomqa}

\section{退行、视界与模型的极限}

哈勃定律将退行速率
\begin{equation}
v_{\rm rec}=HD
\label{eq:c02-recession}
\end{equation}
赋予在选定的宇宙时切片上处于固有距离\(D\)处的物体。对于足够大的\(D\)，方程~\eqref{eq:c02-recession}可以超过\(c\)。这并不违反狭义相对论，因为\(v_{\rm rec}\)并不是在单个局部惯性系中测得的邻近物体的速度。它比较的是变化几何中相距遥远的共动世界线。每个局部观测者仍然测得光速为\(c\)，并且没有任何局部物质物体以超光速越过另一个物体。

\begin{classroomqa}
\audiencequestion{足够遥远的星系能否以超光速退行？}
\lecturerresponse{可以，其宇宙学退行速率\(HD\)可以超过\(c\)。此时没有任何局部物体以超光速越过邻近的观测者，这一结果也不能用于超光速信号传递。它陈述的是遥远共动世界线之间固有距离的变化，而不是在单个事件处测得的普通狭义相对论速度。\lecturetimestamp{01:30:17--01:32:13}}
\end{classroomqa}

超光速退行本身并不能决定现在发出的光是否终将到达我们。这是一个关于整个未来膨胀历史的问题。在减速的物质主导模型中，目前退行快于光的区域以后可以变为可观测的，因为相关视界会演化。如果加速膨胀持续下去，事件视界可能永久地隐藏足够遥远的事件。

\begin{classroomqa}
\audiencequestion{一旦星系的退行速率超过光速，它们就会消失吗？}
\lecturerresponse{并非如此简单。穿越哈勃半径与真正的事件视界是不同的概念。在减速模型中，来自目前超光速退行区域的光仍可能稍后到达我们。持续的暗能量主导会改变整体因果结构，并可能产生一个终极视界。要判定哪些信号会到达，需要时空几何，而不是\(HD\)的瞬时快照。\lecturetimestamp{01:32:14--01:34:18}}
\end{classroomqa}

狭义相对论的语言在全局应用时也可能产生误导。粒子的不变静止质量并不随速度增加。其能量和动量取决于观测者的局部参考系。有质量粒子不能被加速到超过局部光速，但遥远共动星系的退行并不是它穿过单个局部参考系。

\begin{classroomqa}
\audiencequestion{快速退行的星系会获得更大的相对论质量吗？}
\lecturerresponse{最好不要使用``相对论质量''。该星系具有不变静止质量，而其能量和动量依赖于参考系。此外，跨越宇宙学距离赋予的退行速率并不是狭义相对论公式中出现的局部速度。只有用时空度规比较遥远事件之后，这些概念才一致。\lecturetimestamp{01:38:34--01:40:05}}
\end{classroomqa}

\begin{classroomqa}
\audiencequestion{当遥远的退行速率为超光速时，时间和空间轴会发生什么？}
\lecturerresponse{并不存在一个单一的全局洛伦兹参考系，其倾斜的轴可以解决这个问题。狭义相对论比较的是平直时空中的惯性系；宇宙学则需要一个在延展的弯曲时空中联系局部参考系的度规。牛顿式推理在小的局部区域内仍然可靠，但全局构造必须等待广义相对论。\lecturetimestamp{01:40:07--01:42:30}}
\end{classroomqa}

牛顿能量守恒已经为物质给出了正确的标度因子方程，甚至预见了曲率项。波动力学和热力学提供了辐射的标度关系。然而，空间几何、全局能量和视界并不是局部的牛顿问题；它们需要接下来给出的度规描述。

\section{物质、辐射与曲率}

仅含物质的临界解满足
\begin{equation}
\rho_m\propto a^{-3},
\qquad
a(t)\propto t^{2/3}.
\end{equation}
恢复单位质量的牛顿能量给出
\begin{equation}
H^2=\frac{8\pi G}{3}\rho_m+\frac{C}{a^2}.
\end{equation}
正的\(C\)给出永恒膨胀，晚期\(a\sim t\)，\(C=0\)给出临界逃逸，负的\(C\)给出最大标度因子随后再坍缩。定义\(k=-C\)将同一项置于相对论弗里德曼方程的几何一侧。

辐射与物质不同，因为膨胀既改变数密度，也改变每个量子的能量：
\begin{equation}
\begin{aligned}
\lambda&\propto a,
&\qquad E_\gamma&\propto a^{-1},\\
\rho_r&\propto a^{-4},
& a(t)&\propto t^{1/2}.
\end{aligned}
\end{equation}
因此，双组分方程
\begin{equation}
H^2=\frac{C_m}{a^3}+\frac{C_r}{a^4}
\end{equation}
描述了一个辐射主导的早期宇宙，随后是一个物质主导的时代。暗能量、空间几何以及视界的因果结构仍处于这个模型之外，而这正是牛顿近似告诉我们应当把问题交给广义相对论的地方。
